%==============================================================================
% Chapter 5: Proposed Methodology and Implementation
%==============================================================================
\chapter{Proposed Methodology and Implementation}
\label{ch:implementation}

\section{End-to-End Pipeline Overview}
\label{sec:impl_pipeline}

The core methodology transforms raw security telemetry into a validated, actionable incident advisory through eleven sequential stages executed by the \texttt{SOCAnalysisOrchestrator}. Figure \ref{fig:methodology_flow} illustrates the end-to-end dataflow pipeline spanning Telemetry Ingestion, Entity-Anchored Dual Retrieval, and Sovereign Citation Governance.

\begin{figure}[!htbp]
\centering
\resizebox{0.96\textwidth}{!}{%
\begin{tikzpicture}[
    stage/.style={
        rectangle, rounded corners=5pt, draw=teal!70!black, top color=teal!10, bottom color=teal!5,
        very thick, text width=2.6cm, minimum height=1.25cm, align=center, font=\footnotesize\sffamily,
        drop shadow={opacity=0.10, shadow xshift=1pt, shadow yshift=-1pt},
        execute at begin node={\hyphenpenalty=10000\exhyphenpenalty=10000}
    },
    retrieval/.style={
        rectangle, rounded corners=5pt, draw=blue!70!black, top color=blue!10, bottom color=blue!5,
        very thick, text width=2.6cm, minimum height=1.25cm, align=center, font=\footnotesize\sffamily,
        drop shadow={opacity=0.10, shadow xshift=1pt, shadow yshift=-1pt},
        execute at begin node={\hyphenpenalty=10000\exhyphenpenalty=10000}
    },
    reasoning/.style={
        rectangle, rounded corners=5pt, draw=purple!70!black, top color=purple!10, bottom color=purple!5,
        very thick, text width=2.6cm, minimum height=1.25cm, align=center, font=\footnotesize\sffamily,
        drop shadow={opacity=0.10, shadow xshift=1pt, shadow yshift=-1pt},
        execute at begin node={\hyphenpenalty=10000\exhyphenpenalty=10000}
    },
    arrow/.style={-{Stealth[length=6.5pt, width=4.5pt]}, very thick, draw=black!75},
    lbl/.style={font=\scriptsize\bfseries\sffamily, above=3pt}
]

    % Row 1: Telemetry Ingestion (Left to Right)
    \node[stage] (s1) at (0, 0) {\textbf{1. Ingestion \& Schema}\\[2pt]\scriptsize Log Normalization};
    \node[stage] (s2) at (5.2, 0) {\textbf{2. Observable Mining}\\[2pt]\scriptsize Context Extractor};
    \node[stage] (s3) at (10.4, 0) {\textbf{3. Query Formulation}\\[2pt]\scriptsize Search Payload};

    % Row 2: Dual Knowledge Retrieval (Right to Left)
    \node[retrieval] (s4) at (10.4, -3.0) {\textbf{4. Dual-Pass Retrieval}\\[2pt]\scriptsize Bitset + HNSW};
    \node[retrieval] (s5) at (5.2, -3.0) {\textbf{5. Evidence Grounding}\\[2pt]\scriptsize Score Ranking};
    \node[retrieval] (s6) at (0, -3.0) {\textbf{6. Prompt Formulation}\\[2pt]\scriptsize Anti-Injection};

    % Row 3: Sovereign Inference & Governance (Left to Right)
    \node[reasoning] (s7) at (0, -6.0) {\textbf{7. Sovereign LLM}\\[2pt]\scriptsize Qwen3:8B ($T=0$)};
    \node[reasoning] (s8) at (5.2, -6.0) {\textbf{8. Citation Validation}\\[2pt]\scriptsize Reference Audit};
    \node[reasoning] (s9) at (10.4, -6.0) {\textbf{9. Persistence \& Trace}\\[2pt]\scriptsize PostgreSQL \& SSE};

    % Horizontal Connectors with generous margins (>5mm clearance on both sides)
    \draw[arrow] (s1.east) -- (s2.west);
    \node[lbl, color=teal!85!black] at ($(s1.east)!0.5!(s2.west)$) {Alert JSON};

    \draw[arrow] (s2.east) -- (s3.west);
    \node[lbl, color=teal!85!black] at ($(s2.east)!0.5!(s3.west)$) {IoC Set};

    \draw[arrow] (s4.west) -- (s5.east);
    \node[lbl, color=blue!85!black] at ($(s4.west)!0.5!(s5.east)$) {Candidates};

    \draw[arrow] (s5.west) -- (s6.east);
    \node[lbl, color=blue!85!black] at ($(s5.west)!0.5!(s6.east)$) {Top Chunks};

    \draw[arrow] (s7.east) -- (s8.west);
    \node[lbl, color=purple!85!black] at ($(s7.east)!0.5!(s8.west)$) {Advisory};

    \draw[arrow] (s8.east) -- (s9.west);
    \node[lbl, color=purple!85!black] at ($(s8.east)!0.5!(s9.west)$) {Audit Trace};

    % Vertical Connectors with pill badges centered in open gap
    \draw[arrow] (s3.south) -- (s4.north);
    \node[font=\scriptsize\bfseries\sffamily\color{blue!90!black}, fill=white, draw=blue!40, rounded corners=2pt, inner sep=2pt] at ($(s3.south)!0.5!(s4.north)$) {Search Payload};

    \draw[arrow] (s6.south) -- (s7.north);
    \node[font=\scriptsize\bfseries\sffamily\color{purple!90!black}, fill=white, draw=purple!40, rounded corners=2pt, inner sep=2pt] at ($(s6.south)!0.5!(s7.north)$) {Prompt Context};

    % Layer Background Bounding Boxes (inner ysep=6pt leaves 1.6cm open gap between rows)
    \begin{scope}[on background layer]
        \node[fill=teal!4, draw=teal!35, rounded corners=6pt, fit=(s1) (s3), inner xsep=8pt, inner ysep=8pt] (layer1) {};
        \node[fill=blue!4, draw=blue!35, rounded corners=6pt, fit=(s6) (s4), inner xsep=8pt, inner ysep=8pt] (layer2) {};
        \node[fill=purple!4, draw=purple!35, rounded corners=6pt, fit=(s7) (s9), inner xsep=8pt, inner ysep=8pt] (layer3) {};
    \end{scope}

    % Clean Layer Header Badges
    \node[font=\tiny\bfseries\sffamily\color{teal!80!black}, fill=white, draw=teal!35, rounded corners=3pt, inner sep=2.5pt] at (layer1.north) {LAYER 1: TELEMETRY INGESTION \& PREPROCESSING};
    \node[font=\tiny\bfseries\sffamily\color{blue!80!black}, fill=white, draw=blue!35, rounded corners=3pt, inner sep=2.5pt] at (layer2.north) {LAYER 2: ENTITY-ANCHORED DUAL KNOWLEDGE RETRIEVAL};
    \node[font=\tiny\bfseries\sffamily\color{purple!80!black}, fill=white, draw=purple!35, rounded corners=3pt, inner sep=2.5pt] at (layer3.north) {LAYER 3: SOVEREIGN REASONING, CITATION GOVERNANCE \& PERSISTENCE};

\end{tikzpicture}%
}
\caption{Algorithmic pipeline workflow of the AI-Powered SOC Assistant, illustrating the 9 sequential stages across Ingestion, Entity-Anchored Retrieval, and Citation-Governed Sovereign Reasoning.}
\label{fig:methodology_flow}
\end{figure}

\section{Alert Ingestion and Security Normalization}
\label{sec:impl_normalization}
When an alert is received, \texttt{AlertService} maps the incoming JSON or Syslog string into the canonical \texttt{NormalizedSecurityAlert} model:
\begin{itemize}
    \item Resolves timestamps into standardized UTC ISO-8601 strings.
    \item Cleans IPv4/IPv6 strings, stripping port numbers and whitespace.
    \item Maps numerical severities ($0 \le s \le 16$) to ratings (\texttt{LOW}, \texttt{MEDIUM}, \texttt{HIGH}, \texttt{CRITICAL}).
    \item Parses MITRE technique identifiers from Wazuh XML metadata (\texttt{rule.mitre.id}).
\end{itemize}

\section{Deterministic Security Context Extraction}
\label{sec:impl_context}
The second stage executes observable mining via \texttt{ContextExtractor}. Operating deterministically without probabilistic LLM calls, it structures:
\begin{enumerate}[itemsep=1pt,topsep=2pt]
    \item \textbf{Network Indicators:} Source and destination IP addresses and ports.
    \item \textbf{Host Attributes:} Usernames, agent hostnames, operating systems, and process binary paths.
    \item \textbf{Keyword Tokenization:} Tokenizes rule descriptions via regex, filtering against stop-words.
    \item \textbf{Attack Term Mining:} Matches tokens against a taxonomy of 38 attack signals (e.g., \texttt{brute}, \texttt{injection}, \texttt{ransomware}, \texttt{escalation}, \texttt{exfiltration}).
\end{enumerate}

\begin{lstlisting}[language=Python, caption={Deterministic keyword filtering in \texttt{ContextExtractor}.}]
def _keywords_from_text(text: Optional[str]) -> List[str]:
    if not text:
        return []
    tokens = [t.lower().strip(".:,;()[]") for t in _TOKEN_RE.findall(text)]
    return [t for t in tokens if t not in _STOP_WORDS
            and t not in _NOISE_TOKENS and len(t) >= 3]
\end{lstlisting}

\section{Categorized Retrieval Query Construction}
\label{sec:impl_query_builder}
Query compilation in \texttt{RetrievalQueryBuilder} synthesizes up to three distinct query categories:
\begin{enumerate}[itemsep=1pt,topsep=2pt]
    \item \textbf{MITRE Technique Queries:} Created if explicit technique IDs exist (e.g., $Q_{\text{mitre}} = \text{``MITRE ATT\&CK technique } T\text{''}$).
    \item \textbf{NVD / CVE Queries:} Synthesized on CVE regex matches ($Q_{\text{nvd}} = \text{``CVE details''}$).
    \item \textbf{General Security Queries:} Built from normalized event types and top descriptive keywords.
\end{enumerate}

\section{Offline Knowledge Ingestion and Chunking Pipeline}
\label{sec:impl_chunking}
Knowledge ingestion runs offline via \texttt{app/rag/ingest\_knowledge.py}:
\begin{itemize}[itemsep=1pt,topsep=2pt]
    \item \textbf{MITRE ATT\&CK Enterprise Matrix v14.1:} Parsed from STIX 2.1 JSON into 1,618 technique chunks.
    \item \textbf{NIST National Vulnerability Database (NVD):} Indexed into 144 structured CVE chunks.
\end{itemize}
Total indexed corpus size: \textbf{1,762 vector chunks}, embedded into 384-dimensional dense vectors using \texttt{all-MiniLM-\allowbreak L6-v2} alongside metadata payloads (\texttt{source\_id}, \texttt{tactics}, \texttt{title}).

\section{Dual-Pass Entity-Anchored Retrieval Formulation}
\label{sec:impl_retriever}
To eliminate semantic drift, \texttt{Retriever} executes a \textbf{dual-pass search}:
\begin{enumerate}[itemsep=1pt,topsep=2pt]
    \item \textbf{Pass 1 (Payload-Filtered):} For queries carrying MITRE or CVE IDs, Qdrant executes exact payload filtering matching the target identifier.
    \item \textbf{Pass 2 (Dense Semantic Search):} The query vector $\mathbf{v}_q$ is searched across the corpus using unconstrained cosine similarity.
    \item \textbf{Merging and Deduplication:} Results from both passes are merged, retaining the highest genuine cosine score per chunk.
\end{enumerate}

\begin{lstlisting}[language=Python, caption={Dual-pass hybrid search implementation in \texttt{app/rag/retriever.py}.}]
for payload_filter in search_passes:
    pass_results = self._index.search(
        query_vector=vector, top_k=top_k, payload_filter=payload_filter
    )
    for res in pass_results:
        cid = res.chunk_id
        if cid not in accumulated:
            accumulated[cid] = {"chunk": res, "best_score": res.similarity_score}
        elif res.similarity_score > accumulated[cid]["best_score"]:
            accumulated[cid]["best_score"] = res.similarity_score
\end{lstlisting}

\section{Tri-Mode Analysis Execution Paradigms}
\label{sec:impl_tri_mode}
To scientifically evaluate retrieval impact, the orchestrator implements three operational modes:
\begin{enumerate}
    \item \textbf{Mode 1: LLM-Only (\texttt{llm\_only}):} Retrieval is bypassed. The raw alert JSON is evaluated directly by Qwen3:8B, measuring pure parametric pre-training memorization.
    \item \textbf{Mode 2: Standard RAG (\texttt{standard\_rag}):} Unanchored dense cosine similarity over the alert description; top-5 nearest chunks are injected into context.
    \item \textbf{Mode 3: Evidence-Aware RAG (\texttt{evidence\_aware}):} Our contribution: observables are extracted, dual-pass retrieval executes, and citation validation is enforced.
\end{enumerate}

\section{Prompt Formulation and Anti-Injection Constraints}
\label{sec:impl_prompts}
The orchestrator compiles a system prompt enforcing defensive constraints and structured JSON:
\begin{lstlisting}[caption={Standardized JSON output schema required from the LLM.}]
{
  "summary": "1-2 sentence executive summary",
  "severity_assessment": "Technical rationale for assigned severity",
  "observed_indicators": ["192.168.1.105", "admin", "T1110"],
  "suspected_techniques": ["T1110.001"],
  "reasoning": "Detailed forensic explanation",
  "investigation_steps": ["Check auth logs", "Review firewall"],
  "confidence": "HIGH", "limitations": "Missing packet capture",
  "evidence_references": [{"evidence_id": "Exact matching ID", "claim": "...", "supporting_text": "..."}]
}
\end{lstlisting}
Rules mandate: (1) NEVER invent evidence IDs; (2) State missing evidence in limitations; (3) All investigation steps must be non-destructive; (4) Distinguish facts from hypotheses.

\section{Automated Citation Reference Validation Engine}
\label{sec:impl_citation_validator}
To eliminate generative hallucinations, \texttt{\_validate\_evidence()} parses the generated JSON and cross-examines citations against the actually retrieved evidence dictionary:
\begin{itemize}
    \item \textbf{Reference Existence Check:} Valid citation IDs are stored in PostgreSQL with relevance scores.
    \item \textbf{Hallucination Trapping:} Invented IDs are stripped from official evidence, and an explicit warning is appended to limitations: \texttt{Note: LLM hallucinated unknown evidence reference 'X'}.
\end{itemize}

\begin{lstlisting}[language=Python, caption={Citation validation and hallucination trapping in \texttt{orchestrator.py}.}]
def _validate_evidence(self, analysis: SOCAnalysis, evidence_set: Optional[EvidenceSet]):
    valid_dict = {ev.evidence_id: ev for ev in evidence_set.evidence} if evidence_set else {}
    validated_refs, evidence_creates = [], []
    for ref in analysis.evidence_references:
        if ref.evidence_id in valid_dict:
            validated_refs.append(ref)
            ev = valid_dict[ref.evidence_id]
            evidence_creates.append(EvidenceCreate(
                source=ev.source, title=ev.title, relevance_score=ev.similarity_score,
                citation=f"{ev.source_identifier} - {ev.title}"
            ))
        else:
            analysis.limitations += f"\nNote: LLM hallucinated unknown reference '{ref.evidence_id}'."
    return validated_analysis, evidence_creates
\end{lstlisting}

\section{Analyst User Interface and Interactive Entity Graph}
\label{sec:impl_frontend}
The presentation layer is a React 19 single-page console built with TypeScript and Tailwind CSS:
\begin{itemize}
    \item \textbf{Live Triage Queue:} Ingestion queue with severity badges, technique tags, and source filters.
    \item \textbf{D3.js Entity Graph:} Force-directed canvas displaying relationships between alerts, IPs, users, and MITRE techniques.
    \item \textbf{Real-Time Trace Waterfall:} Visualizes pipeline execution streaming over SSE with stage timings.
    \item \textbf{Technical Inspector Drawer:} Displays raw JSON payloads, cosine scores, and citation provenance.
\end{itemize}